\documentclass{sprawozdanie-agh}

\usepackage[utf8]{inputenc}
\usepackage{matlab-prettifier}
\usepackage{siunitx}
\usepackage{booktabs,caption,ragged2e}

\makeatletter

% ========= Początek Dokumentu ========= 
\begin{document}
% ========= Strona tytułowa ========= 
\przedmiot{Teoria Sterowania II: Ćwiczenia Laboratoryjne}
\tytul{Wytyczne dla Ćwiczenia 6}
\podtytul{Optymalizacja Parametryczna}
\kierunek{Automatyke i~Robotyka}
\autor{Dariusz Cieślar, Maciej Różewicz}
\data{Kraków, \today}

\stronatytulowa{}
% ========= Treść Właściwa =========
\section{Cel ćwiczenia}
Celem ćwiczenia jest zapoznanie się ze wskaźnikami jakości układów sterowania oraz optymalizacją parametryczną układów sterowania ze względu na wybrane wskaźniki.

Zadanie optymalizacji parametrycznej dotyczy problemów, gdzie struktura regulatora i~model matematyczny obiektu są stałe, a~optymalizowane są jedynie wartości niektórych parametrów.
Dobór odpowiedniej funkcji celu, która jest optymalizowana, jest zagadnieniem zasadniczym.
Zwykle żąda się, aby układ regulacji możliwie dokładnie i~szybko odtwarzał wartość zadaną, z~drugiej zaś strony, by był możliwie odporny na sygnały zakłócające.
Często wymagania te są sprzeczne i zadaniem projektanta jest znalezienie satysfakcjonującego kompromisu. 
Znaczącą zaletą procedury optymalizacji parametrycznej jest ustalenie pożądanych cech na wstępnym etapie projektowania układu sterowania.
Procedura ta jest zwykle iteracyjna: funkcja celu ewoluuje w trakcie procesu projektowania, jednak na każdym etapie możemy być pewni, że dla sfomułowanej funkcji celu uzyskujemy najlepsze możliwe nastawy regulatora.

%%%%%%%%%%%%%%%%%%%%%%%%%%%%%%%%%%%%%%%%%%%%%%%%%%%%%%%%%%%
\section{Przebieg ćwiczenia}

\subsection{Wyznaczanie wskaźników jakości}

\begin{itemize}
    \item serwo
\end{itemize}

    %======================================================
    \subsection{Optymalizacja parametryczna: wskaźniki całkowe}
    Dla obiektu opisywanego transmitancją:
    \begin{equation}
        \label{eq:system_1}
        G_{O}(s) = \frac{2s+3}{s^3+2s^2+s},
    \end{equation}
    w~układzie regulacji automatycznej, rozważyć do optymalizacji parametrycznej regulatory:
    \begin{itemize}
        \item P: $G_{R,P}(s) = K_{P}$,
        \item PD: $G_{R,PD}(s) = K_{P}(1 + T_{D}s)$.
    \end{itemize}

    \begin{description}
    \item[Krok 5] utworzyć model symulacyjny układu sterowania w~Simulinku (łącznie z~wejściami zakłóceń),
    \item[Krok 6] wyznaczyć analitycznie ograniczenia na zakres zmiennych projektowych (w zależności od zadanej struktury regulatora) - zaznaczyć ten obszar na płaszczyźnie,
    \item[Krok 7] wykonać numerycznie optymalizacje ze względu na całkę z kwadratu uchybu sterowania (ang. \textit{Integral
            Square Error} - \textbf{ISE}):
                \begin{equation}
                    \label{eq:j1}
                    J_{ISE} = \int_{0}^{\infty}{e(t)^{2}dt},
                \end{equation}
    \item[Krok 8] wykonać numerycznie optymalizację ze względu na całkę z kwadratu wyjścia układu regulacji (ang.
            \textit{Integral Square Output} - \textbf{ISO}) dla stabilizacji w zerze:
                \begin{equation}
                    \label{eq:j7}
                    J_{ISO} = \int_{0}^{\infty}{y^{2}(t)dt}
                \end{equation}
    \item[Krok 9] wykonać numerycznie optymalizację ze względu na odporność na obciążenie,
    \item[Krok 10] wykonać numerycznie optymalizację ze względu na wskaźniki \textbf{ISE} i~na obciążenie,
    \item[Krok 11] dla każdego zestawu dobranych nastaw regulatora wykonać symulacyjnie następujące eksperymenty i pokazać wartości sterowania na jednym rysunku:
        \begin{itemize}
            \item odpowiedź na skok jednostkowy dla wartości zadanej, bez zakłóceń,
            \item odpowiedź na skok jednostkowy z~zakłóceniem obciążeniowym (skok jednostkowy, ale w~innej chwili niż wartość zadana).
        \end{itemize}
    \end{description}

    %======================================================
    \subsection{Optymalizacja parametryczna: wskaźniki odpowiedzi czasowej}

    \begin{description}
    \item[Krok 12] Wyznaczyć optymalny regulator proporcjonalny $G_{R}(s) = K$, dla obiektu opisywanego transmitancją:
    $$
        G_{O}(s) = \frac{1}{s^2 + s + 1},
    $$
    minimalizujący czas narastania, przy przeregulowaniu nie przekraczającym $10\%$.
    \end{description}

\end{document}
